% ---
% id: receiver/component-catalogue/rem:entry-roles
% kind: rem
% title: "Physical roles of the catalogue entries"
% summary: >-
%   The mirror's handedness flip is a choice of frame, the phase shifter is the
%   interferometer's control knob, and fibre, coupling, filter and connector
%   losses are loss entries at different sites. Route (ii) is the special case
%   of route (iii) with one generic block spanning the receiver.
% prominence: standard
% status: accepted
% uses: [thm:catalogue-admissibility]
% mentions: [receiver/detector-imperfections/thm:efficiency, def:catalogue-entries]
% assumes: []
% refs: []
% provenance: newly-derived
% lean:
% ---

\begin{remark}\label{rem:entry-roles}
Read physically, the entries of Definition~\ref{def:catalogue-entries} play
the following roles.
\begin{itemize}[leftmargin=1.6em]
\item \emph{C2.} Entry C2 is the polarization-dependent beamsplitter of
  entry C1 with its per-output losses, written in the mode order
  $(p_1H,p_1V,p_2H,p_2V)$; its ideal preset $t_H=1$, $t_V=0$ is the extreme
  case, the polarizing beamsplitter.
\item \emph{C3.} The preset $\mathrm{HWP}(\pi/8)$ is the $22.5^{\circ}$
  plate that maps the rectilinear basis to the diagonal one.
\item \emph{C5.} The handedness flip is a declared choice of coordinate frame,
  not a physical claim.
\item \emph{C6.} This is the interferometric receiver's control knob.
\item \emph{C7.} This is the loss primitive: fibre attenuation, coupling
  efficiency, spectral-filter insertion loss and connector loss are all C7 at
  different \emph{sites}, and the sites are what Theorem~\ref{thm:efficiency} is
  about.
\item \emph{C8.} It is contractive and not unitary unless $t=t_\perp=1$: the
  blocked polarization is loss.
\item \emph{C10.} This entry makes route (ii) a special case of route (iii).
  A user with a measured \tref{instrument matrix}, or with a component the
  catalogue does not have, declares one block, and a route-(ii) declaration
  is recovered as the route-(iii) declaration of exactly one C10 block
  spanning the whole receiver.
\end{itemize}
\end{remark}
